\documentclass[10pt,a4paper]{amsart}
\usepackage{amsmath,amssymb,amsthm,graphicx}
\usepackage{verbatim}
\newenvironment{reference}{\comment}{\endcomment}
%\newenvironment{reference}{}{}
\newenvironment{slogan}{\comment}{\endcomment}
\newenvironment{history}{\comment}{\endcomment}

% For commutative diagrams we use Xy-pic
\usepackage[all]{xy}

% We use 2cell for 2-commutative diagrams.
\xyoption{2cell}
% Theorem environments.
%
\theoremstyle{plain}
\newtheorem{theorem}[subsection]{Theorem}
\newtheorem{proposition}[subsection]{Proposition}
\newtheorem{lemma}[subsection]{Lemma}

\theoremstyle{definition}
\newtheorem{definition}[subsection]{Definition}
\newtheorem{example}[subsection]{Example}
\newtheorem{exercise}[subsection]{Exercise}
\newtheorem{situation}[subsection]{Situation}

\theoremstyle{remark}
\newtheorem{remark}[subsection]{Remark}
\newtheorem{remarks}[subsection]{Remarks}

\numberwithin{equation}{subsection}

% Macros
%
\def\lim{\mathop{\mathrm{lim}}\nolimits}
\def\colim{\mathop{\mathrm{colim}}\nolimits}
\def\Spec{\mathop{\mathrm{Spec}}}
\def\Hom{\mathop{\mathrm{Hom}}\nolimits}
\def\Ext{\mathop{\mathrm{Ext}}\nolimits}
\def\SheafHom{\mathop{\mathcal{H}\!\mathit{om}}\nolimits}
\def\SheafExt{\mathop{\mathcal{E}\!\mathit{xt}}\nolimits}
\def\Sch{\mathit{Sch}}
\def\Mor{\mathop{\mathrm{Mor}}\nolimits}
\def\Ob{\mathop{\mathrm{Ob}}\nolimits}
\def\Sh{\mathop{\mathit{Sh}}\nolimits}
\def\NL{\mathop{N\!L}\nolimits}
\def\CH{\mathop{\mathrm{CH}}\nolimits}
\def\proetale{{pro\text{-}\acute{e}tale}}
\def\etale{{\acute{e}tale}}
\def\QCoh{\mathit{QCoh}}
\def\Ker{\mathop{\mathrm{Ker}}}
\def\Im{\mathop{\mathrm{Im}}}
\def\Coker{\mathop{\mathrm{Coker}}}
\def\Coim{\mathop{\mathrm{Coim}}}

% Boxtimes
%
\DeclareMathSymbol{\boxtimes}{\mathbin}{AMSa}{"02}

%
% Macros for moduli stacks/spaces
%
\def\QCohstack{\mathcal{QC}\!\mathit{oh}}
\def\Cohstack{\mathcal{C}\!\mathit{oh}}
\def\Spacesstack{\mathcal{S}\!\mathit{paces}}
\def\Quotfunctor{\mathrm{Quot}}
\def\Hilbfunctor{\mathrm{Hilb}}
\def\Curvesstack{\mathcal{C}\!\mathit{urves}}
\def\Polarizedstack{\mathcal{P}\!\mathit{olarized}}
\def\Complexesstack{\mathcal{C}\!\mathit{omplexes}}
% \Pic is the operator that assigns to X its picard group, usage \Pic(X)
% \Picardstack_{X/B} denotes the Picard stack of X over B
% \Picardfunctor_{X/B} denotes the Picard functor of X over B
\def\Pic{\mathop{\mathrm{Pic}}\nolimits}
\def\Picardstack{\mathcal{P}\!\mathit{ic}}
\def\Picardfunctor{\mathrm{Pic}}
\def\Deformationcategory{\mathcal{D}\!\mathit{ef}}
\usepackage{cleveref}
\crefname{originallemma}{original lemma}{original lemmas}
\Crefname{originallemma}{Original lemma}{Original lemmas}
\makeatletter
\newlabel{duality-lemma-dualizing-module-proper-over-A}{{0AWP}{source}{}{sourcepointer-1}{}}
\newlabel{duality-lemma-dualizing-module-proper-over-A@cref}{{[originallemma][0][]0AWP}{[1][0][]source}{}{}{}}
\makeatother
\crefname{originalreference}{original reference}{original references}
\Crefname{originalreference}{Original reference}{Original references}
\makeatletter
\newlabel{lemma-globally-generated}{{01AN}{source}{}{sourcepointer-2}{}}
\newlabel{lemma-globally-generated@cref}{{[originalreference][0][]01AN}{[1][0][]source}{}{}{}}
\makeatother
\crefname{originalreference}{original reference}{original references}
\Crefname{originalreference}{Original reference}{Original references}
\makeatletter
\newlabel{situation-Cohen-Macaulay-curve}{{0B5D}{source}{}{sourcepointer-3}{}}
\newlabel{situation-Cohen-Macaulay-curve@cref}{{[originalreference][0][]0B5D}{[1][0][]source}{}{}{}}
\makeatother
\title{An extension of the dualizing sheaf by a zero-dimensional quotient with at least two sections is globally generated}
\author{The Stacks Project Authors}
\date{stacks-project:curves.tex:local-snapshot-20261012; GNU FDL 1.2 or later}
\begin{document}
\maketitle
\noindent\textit{Editorial note.} All delivered mathematical text is the author's own, copied from the pinned Stacks Project chapter file curves.tex; no statement or proof step was rewritten, reordered or removed.\par\medskip
\noindent\textit{Editorial note.} The counted proof is the author's complete argument: the rank-1 setup, the H\textasciicircum{}1 duality step, the self-map closure argument, the author's marked automorphism fact used as a step, and the final contradiction are all retained verbatim.\par\medskip
\noindent\textit{Editorial note.} The label lemma-globally-generated-curve names the counted statement, and the proof environment immediately following it in the source is delivered verbatim as its proof.\par\medskip
\noindent\textit{Editorial note.} The counted statement is carried by the author's own lemma environment in the chapter source.\par\medskip
\noindent\textit{Allowed standard prerequisites (editorial).}\begin{itemize}
\item The globally generated subsheaf lemma. Same chapter, Lemma lemma-globally-generated. Shows the subsheaf F' generated by global sections already surjects onto Q, so F' is nonzero.
\item Duality on proper curves. Duality for Schemes, Lemma duality-lemma-dualizing-module-proper-over-A. Turns nonzero H\textasciicircum{}1(X, F') into a nonzero map F' -> omega\_X in the contradiction step.
\end{itemize}\medskip
\ifdefined\hypertarget\hypertarget{sourcepointer-1}{}\fi
\noindent\textit{Source reference.} Original lemma 0AWP: sources/stacks-chow/tags.txt, source lines 13727--13727. Prerequisite pointer invoked at proof line(s) 1426; it names a standard fact of a sibling chapter that the counted argument uses, and every step of the counted proof built on it is retained verbatim in the delivered body.\par\medskip
\ifdefined\hypertarget\hypertarget{sourcepointer-2}{}\fi
\noindent\textit{Source reference.} Original reference 01AN: sources/stacks-curves/pinned-original-source.tex, source lines 1372--1372. Prerequisite pointer invoked at proof line(s) 1418; it names a standard fact of the same chapter that the counted argument uses, and every step of the counted proof built on it is retained verbatim in the delivered body.\par\medskip
\ifdefined\hypertarget\hypertarget{sourcepointer-3}{}\fi
\noindent\textit{Source reference.} Original reference 0B5D: sources/stacks-curves/pinned-original-source.tex, source lines 1307--1307. Prerequisite pointer invoked at proof line(s) 1408; it names a standard fact of the same chapter that the counted argument uses, and every step of the counted proof built on it is retained verbatim in the delivered body.\par\medskip
\begin{lemma}
\label{lemma-globally-generated-curve}
In Situation \ref{situation-Cohen-Macaulay-curve} assume that
$X$ is integral. Let $0 \to \omega_X \to \mathcal{F} \to \mathcal{Q} \to 0$
be a short exact sequence of coherent $\mathcal{O}_X$-modules with
$\mathcal{F}$ torsion free, $\dim(\text{Supp}(\mathcal{Q})) = 0$,
and $\dim_k H^0(X, \mathcal{Q}) \geq 2$. Then $\mathcal{F}$
is globally generated.
\end{lemma}

\begin{proof}
Consider the submodule $\mathcal{F}'$ generated by the global sections. By
Lemma \ref{lemma-globally-generated} we see that $\mathcal{F}' \to \mathcal{Q}$
is surjective, in particular $\mathcal{F}' \not = 0$. Since $X$ is a curve, we
see that $\mathcal{F}' \subset \mathcal{F}$ is an inclusion of rank $1$
sheaves, hence $\mathcal{Q}' = \mathcal{F}/\mathcal{F}'$ is supported in
finitely many points. To get a contradiction, assume that
$\mathcal{Q}'$ is nonzero. Then we see that $H^1(X, \mathcal{F}') \not = 0$.
Then we get a nonzero map $\mathcal{F}' \to \omega_X$ by the universal
property (Duality for Schemes, Lemma
\ref{duality-lemma-dualizing-module-proper-over-A}).
The image of the composition $\mathcal{F}' \to \omega_X \to \mathcal{F}$
is generated by global sections, hence is inside of $\mathcal{F}'$.
Thus we get a nonzero self map $\mathcal{F}' \to \mathcal{F}'$.
Since $\mathcal{F}'$ is torsion free of rank $1$ on a proper curve
this has to be an automorphism (details omitted). But then this implies that
$\mathcal{F}'$ is contained in $\omega_X \subset \mathcal{F}$
contradicting the surjectivity of $\mathcal{F}' \to \mathcal{Q}$.
\end{proof}

\end{document}
